\documentclass[11pt]{article}
\usepackage[margin=1.05in]{geometry}
\usepackage{amsmath,amssymb,amsthm,mathtools}
\usepackage{xcolor}
\usepackage[colorlinks,linkcolor=blue!55!black,citecolor=blue!55!black]{hyperref}

\newtheorem{theorem}{Theorem}
\newtheorem{proposition}[theorem]{Proposition}
\newtheorem{lemma}[theorem]{Lemma}
\newtheorem{corollary}[theorem]{Corollary}
\theoremstyle{definition}
\newtheorem{remark}[theorem]{Remark}
\newtheorem{convention}[theorem]{Convention}

\newcommand{\hex}{\mathsf{hex}}
\newcommand{\hexR}{\mathsf{hex}^{\mathsf{R}}}
\newcommand{\trap}{\mathsf{trap}}
\newcommand{\Ktile}{\mathsf{K}}
\newcommand{\tritile}{\mathsf{tri}}
\newcommand{\la}{\lambda}
\newcommand{\Star}{\operatorname{St}}
\newcommand{\cc}{c^{\nu}_{\la\mu}}
\newcommand{\dd}{d^{\nu}_{\la\mu}}

\title{The Pylyavskyy--Yang hexagon is the vertex star\\
\large Q569: the hexagonal puzzle piece is a \emph{bulk} tile,
and its boundary piece is a \emph{truncation}, not a K-matrix}
\author{Clio}
\date{23 October 2026}

\begin{document}
\maketitle

\begin{center}\small
\fbox{\begin{minipage}{0.88\textwidth}\small
\textbf{Provenance} (PROTOCOL §2.3).\\
\textbf{Author:} Clio.\quad\textbf{Date:} 23 October 2026 (prove session, cycle 2).\\
\textbf{Recipient:} Robin Langer (\texttt{langer.robin@gmail.com}).\\
\textbf{Repository:} \texttt{clio-vega/proofs}.\quad
\textbf{Commit:} named in the covering email; this artifact and the scripts in
\texttt{proofs/code-1023c2-Q569/} are committed together.\\
\textbf{Source read:} \texttt{library/2026-10-23-browse-arxiv-raw/ep-1801.07667/}
(\texttt{ktiles.tex}, 1725 lines; \texttt{hex-tiling.tex}, 65 lines), e-print of
Pylyavskyy--Yang, arXiv:1801.07667. Extraction level: \texttt{deep-read}
(tile macros, all four tiling theorems, \S\texttt{ss:K-coefficients},
\S\texttt{ss:protrude} read verbatim in the source).\\
\textbf{Status:} the geometric theorems are \texttt{proved}; the
\S\texttt{ss:protrude} equivalence is \texttt{computed} (census domain stated
in \S\ref{sec:protrude}); the integrability \emph{negatives} are
\texttt{proved} where stated as such and flagged otherwise.
\end{minipage}}
\end{center}

\begin{abstract}
Pylyavskyy and Yang (arXiv:1801.07667) introduce two hexagonal puzzle pieces
$\hex,\hexR$ giving Littlewood--Richardson rules for the $K$-homology of
Grassmannians. Their paper contains no integrability claim; the question Q569
asked whether $\hex$ is a \emph{bulk} $R$-matrix piece or a \emph{boundary}
K-matrix piece satisfying a reflection equation, the latter suggested by
\S\texttt{ss:protrude}, where a distinguished trapezoidal tile is introduced and
explicitly forbidden outside the left edge of the puzzle.

The answer is \emph{neither}, and the reason is a single geometric
identification: \textbf{$\hex$ and $\hexR$ are the two orientation-keyed
alternating labellings of the closed vertex star of the triangular lattice}
(Theorem~\ref{thm:star}) --- the union of the six unit triangles containing a
lattice vertex. Everything the paper observes follows. The hexagon is a bulk
piece indexed by its centre; it fails to fit inside a wall exactly when its
centre \emph{lies on} the wall, and then exactly half of it is outside
(Theorem~\ref{thm:protrude}), which is why enlarging the ambient rectangle by
one suffices; the \S\texttt{ss:protrude} trapezoid is \emph{literally} the
hexagon intersected with the half-plane, with every retained edge label
unchanged (Theorem~\ref{thm:trunc}); and $\hexR$ is the image of $\hex$ under
every element of $D_6$ that exchanges the two triangle orientations
(Corollary~\ref{cor:dual}), so the two theorems are one theorem conjugated by
Grassmannian duality, not two boundaries of one model.

The reflection-equation reading is refuted on three independent grounds
(\S\ref{sec:refl}): the boundary piece has \emph{zero moduli}, being determined
by the bulk piece and the wall; the bulk tile set is \emph{chiral}, so no
lattice reflection preserves it and there is nothing to reflect in; and the
third edge label~$2$ is a deterministic wall-adjacency indicator rather than an
amplitude. Separately, the tiles are not vertices of a line-conserving model at
all (\S\ref{sec:nomodel}): ten of the seventeen carry exactly one $1$-edge in
each of the three direction classes.

All of this rests on a tile extraction verified by four independent positive
controls (\S\ref{sec:controls}): an engine built from set-valued tableaux
reproduces the puzzle counts for all four of the paper's theorems
($208/208$, $208/208$, $187/187$, $54/54$ at $n=5,k=2$; $875/875$ at
$n=6,k=3$). I also verify the paper's \emph{unproved} \S\texttt{ss:protrude}
claim ($187/187$ and $875/875$), and quantify the protrusion loss on the
paper's own example: $d=4$, the unfixed count is $2$, and both of the paper's
fixes give $4$.
\end{abstract}

\section{The question, and the answer}

\textbf{Q569, as posed.} Is the $\hex/\hexR$ tile set a \emph{bulk} $R$-matrix
satisfying the Yang--Baxter equation, or a bulk $R$-matrix \emph{plus} a
one-sided boundary K-matrix satisfying the reflection (Sklyanin) equation?

The evidence for the second reading was three sentences of
\S\texttt{ss:protrude}: ``the $\hex$ piece will protrude to the left of the
puzzle with side length~$7$''; ``(By the way things are set up, the hexagon
never needs to protrude to the right or below.)''; and, of the alternative fix
by an extra trapezoidal piece, ``However, we do not want this piece used
elsewhere.'' A tile that acts at one edge and is forbidden in the interior does
look like a K-matrix.

\textbf{Answer.} Neither. The hexagon is a bulk piece. The trapezoid is a
boundary piece, but it is the bulk piece \emph{clipped by the wall}, and a
clipping is not a reflection: it has no spectral parameter, no amplitude and no
moduli. The one-sidedness that motivated the question is real, and
Theorem~\ref{thm:protrude} explains it exactly --- but it is a statement about
where a vertex star can sit relative to a lattice line, not about an
integrable boundary condition.

\section{Objects: the tiles, extracted}\label{sec:objects}

Everything below is extracted from the e-print's own \texttt{tikz} macros
(\texttt{ktiles.tex}, lines 232--236):
\[
  \verb|\r{i}| = \text{unit edge, direction } 60i^\circ,\ \text{label } 0;\quad
  \verb|\b{i}| = \text{label } 1;\quad
  \verb|\bb{i}| = \text{label } 2 .
\]
Since $\verb|\xcoord|=\cos 60^\circ$ and $\verb|\ycoord|=\sin 60^\circ$
(\texttt{ktiles.tex}, lines 4--5), the paper's \texttt{tikz} coordinates
\emph{are} the triangular-lattice basis. I therefore work throughout in integer
coordinates $(a,b) \leftrightarrow a e_0 + b e_1$ with
$e_0=(1,0)$, $e_1=(\tfrac12,\tfrac{\sqrt3}{2})$, so that
\[
 D_0,\dots,D_5 \;=\; (1,0),\,(0,1),\,(-1,1),\,(-1,0),\,(0,-1),\,(1,-1) .
\]
The shear to Euclidean coordinates is affine, so incidence, convexity and
point-in-polygon are unaffected.

A tile is a closed counterclockwise walk of unit steps carrying labels. The
seventeen tiles of the paper are then, verbatim:

\begin{center}\small
\begin{tabular}{llll}
\hline
tile & source & step sequence $(\text{dir},\text{label})$ & $\#\triangle$\\
\hline
$0$-triangles & fig.~LR-tiles & $(0,0)(2,0)(4,0)$;\ $(1,0)(3,0)(5,0)$ & 1\\
$1$-triangles & fig.~LR-tiles & $(0,1)(2,1)(4,1)$;\ $(1,1)(3,1)(5,1)$ & 1\\
rhombi & fig.~LR-tiles & $(0,1)(1,0)(3,1)(4,0)$, $+$ 2 rotations & 2\\
$\Ktile$ (Vakil) & fig.~K-tiles & $(1,1)(1,0)(3,1)(3,0)(5,1)(5,0)$ & 4\\
$\tritile$ (WZ) & fig.~K-tiles & $(0,1)(0,0)(2,1)(2,0)(4,1)(4,0)$ & 4\\
$\hex$ & fig.~K-tiles & $(1,0)(2,1)(3,0)(4,1)(5,0)(0,1)$ & 6\\
$\hexR$ & fig.~K-tiles & $(0,0)(1,1)(2,0)(3,1)(4,0)(5,1)$ & 6\\
$\trap$ & \S\texttt{ss:protrude} & $(1,0)(2,1)(4,0)(4,1)(0,1)$ & 3\\
\hline
\end{tabular}
\end{center}

\noindent together with the five relabelled pieces of \S\texttt{ss:protrude}
(\S\ref{sec:protrude}). All seventeen close and have positive signed area
(verified).

\begin{remark}[the alternation]
The KTW rhombi carry labels $(1,0,1,0)$ around their boundary and $\hex,\hexR$
carry $(0,1,0,1,0,1)$: both are \emph{alternating}. The trapezoid carries
$(0,1,0,1,1)$ --- alternation broken. So the hexagon is the $m=3$ member of the
family whose $m=2$ member is the KTW rhombus, and the trapezoid is not a member
of that family at all. This is the first sign that the trapezoid is a derived
object rather than a new primitive.
\end{remark}

\begin{convention}[boundary, verified against the paper's own figure]\label{conv:bdry}
Let $n\ge k$ and let the ambient rectangle be $k$ rows of length $n-k$. For a
partition $\la$ inside it, $\sigma(\la)\in\{0,1\}^n$ is the lower-right boundary
path read from top right to bottom left, $0$ for horizontal and $1$ for vertical
(\texttt{ktiles.tex}, lines 526--536). The puzzle region is the upright triangle
with vertices $(0,0),(n,0),(0,n)$, and
\[
\begin{aligned}
\text{bottom, } (0,0)\!\to\!(n,0) \text{ in direction } 0 &: \ \sigma(\nu),\\
\text{left, } (0,0)\!\to\!(0,n) \text{ in direction } 1 &: \ \sigma(\la),\\
\text{right, } (0,n)\!\to\!(n,0) \text{ in direction } 5 &: \ \sigma(\mu).
\end{aligned}
\]
I verified all three against figure \texttt{bdry} (\texttt{ktiles.tex}, lines
610--617), which prints the three sides explicitly for
$\la=(2,1,0),\mu=(3,2,0),\nu=(4,3,1)$, $n=9$, $k=3$: exact match on all
$27$ edges.
\end{convention}

\begin{remark}[the wall]\label{rem:wall}
The \emph{left} side of the puzzle region runs from $(0,0)$ to $(0,n)$, i.e.\ in
direction $1$ (equivalently $4$). This is the side past which the hexagon
protrudes, and it is exactly the direction class on which
\S\texttt{ss:protrude} introduces the third label~$2$. I call the line $a=0$
the \emph{wall} $W$.
\end{remark}

\section{Positive controls}\label{sec:controls}

The brief's first instruction was to confirm a well-defined object exists before
testing any equation on it. I built two engines and made them disagree or agree.

\emph{Engine A (puzzles).} Exact cover of the $n^2$ unit triangles of the region
by lattice translates of the tiles, with unit-edge labels matching between
adjacent tiles and against Convention~\ref{conv:bdry}. Edges interior to a tile
are unconstrained.

\emph{Engine B (the ring, independent).} The single stable Grothendieck
polynomial $G_\la=\sum_T(-1)^{|T|-|\la|}x^T$ over semistandard
\emph{set-valued} tableaux (\texttt{ktiles.tex}, lines 711--716), multiplied and
comultiplied as an actual polynomial, then expanded in the $G$-basis by
degree-triangular peeling ($G_\la=s_\la+\text{higher}$). Engine~B touches no
puzzle and no tableau rule from the paper: it goes through the ring. Controls on
Engine~B alone: the minimal-degree part of $G_\la$ is $s_\la$; and $\cc$ at
$|\nu|=|\la|+|\mu|$ reproduces the classical Littlewood--Richardson
coefficients of $s_{21}^2$.

\begin{center}\small
\begin{tabular}{lllr}
\hline
check & domain & result \\
\hline
KTW count $=\cc$, paper's running example & $n{=}9,k{=}3$, $(2,1),(3,2),(4,3,1)$ & $2=2$ \\
$\Ktile$-puzzles $=\cc$ ($K$-theory, Vakil) & $n{=}5,k{=}2$, all $(\la,\mu,\nu)$ & $208/208$ \\
$\tritile$-puzzles $=\tilde c^\nu_{\la\mu}$ (Wheeler--Zinn-Justin) & $n{=}5,k{=}2$, all & $208/208$ \\
$\hex$-puzzles $=\dd$ ($K$-homology) & $n{=}5,k{=}2$, all & $187/187$ \\
$\hexR$-puzzles $=\tilde d^\nu_{\la\mu}$ & $n{=}5,k{=}2$, all & $54/54$ \\
$\hex$-puzzles $=\dd$, one level up & $n{=}6,k{=}3$, all & $875/875$ \\
\hline
\end{tabular}
\end{center}

\noindent
\textbf{Census domain, stated exactly.} ``All $(\la,\mu,\nu)$'' means: every
triple of partitions inside the ambient rectangle prescribed by the relevant
theorem --- $k\times(n-k)$ for $\Ktile,\tritile$; $k\times(n-k-1)$ for $\hex$;
$(k-1)\times(n-k)$ for $\hexR$ --- subject to the theorem's degree inequality
($|\nu|\ge|\la|+|\mu|$ for $\Ktile,\tritile$; $|\nu|\le|\la|+|\mu|$ for
$\hex,\hexR$). The $n=6,k=3$ row is the ``one level up'' test: it moves $k$ off
$2$, where a two-row coincidence could not be detected.

\section{The hexagon is the vertex star}\label{sec:star}

\begin{theorem}[Vertex star]\label{thm:star}
Let $v$ be a vertex of the triangular lattice and let
$\Star(v)=\bigcup\{T : T \text{ a unit triangle},\ v\in T\}$ be its closed star,
the union of the six unit triangles containing $v$. Then:
\begin{enumerate}
\item[(i)] $\Star(v)$ is the convex hexagon with vertex set
$\{v+D_j\}_{j\in\mathbb{Z}/6}$, and its boundary is the closed hexagonal walk
$v+D_0 \to v+D_1 \to \cdots \to v+D_5 \to v+D_0$ of six unit edges.
\item[(ii)] The six triangles of $\Star(v)$ alternate in orientation around
$v$, and the boundary edge contributed by each is the edge of that triangle
opposite $v$.
\item[(iii)] There are exactly two alternating $\{0,1\}$-labellings of this
hexagonal boundary. They are $\hex$ and $\hexR$, and they are distinguished by
triangle orientation: $\hex$ puts the label $1$ on the outer edge of each
\emph{upward} triangle and $0$ on each \emph{downward} one; $\hexR$ reverses
this.
\end{enumerate}
\end{theorem}

\begin{proof}
(i) The unit triangles containing $v$ are exactly
$T_j=\{v,\,v+D_j,\,v+D_{j+1}\}$ for $j\in\mathbb{Z}/6$, since consecutive
$D_j,D_{j+1}$ are at $60^\circ$ and span a unit triangle with $v$, and every
unit triangle containing $v$ arises this way. Their union is the convex hull of
$\{v+D_j\}$, and the edges $[v+D_j,v+D_{j+1}]$ not containing $v$ concatenate to
the stated closed walk.

For the identification with the paper's tiles: summing the step sequences,
$\hex$ and $\hexR$ each close, and in each case the six vertices minus their
centroid are precisely $D_0,\dots,D_5$ (computed exactly: the vertex coordinate
sums are divisible by $6$, so the centroid is a lattice point $v$, and the
vertex set is $v+D$). Hence both tiles equal $\Star(v)$ as regions.

(ii) $T_j$ and $T_{j+1}$ share the edge $[v,v+D_{j+1}]$ and lie on opposite
sides of it, so they have opposite orientations; the parity of $j$ therefore
determines the orientation class. The boundary edge of $T_j$ is
$[v+D_j,v+D_{j+1}]$, which is the edge of $T_j$ not meeting $v$.

(iii) An alternating labelling of a $6$-cycle is determined by its value on one
edge, so there are exactly two. By (ii) the two alternation classes of edges are
precisely the two orientation classes of triangles, so the two labellings are
the two orientation-keyed ones. Reading the paper's step sequences
(\S\ref{sec:objects}) and matching each boundary edge to its triangle gives: in
$\hex$ the upward triangles' outer edges carry $1$ and the downward ones $0$;
in $\hexR$ the reverse (verified directly).
\end{proof}

\begin{corollary}[$\hexR$ is $\hex$ seen in the other orientation]\label{cor:dual}
Let $g\in D_6$ act on the lattice. Then $g(\hex)=\hexR$ if and only if $g$
exchanges the two triangle orientations; otherwise $g(\hex)=\hex$. Explicitly,
the odd rotations $\rho^{1},\rho^{3},\rho^{5}$ and the three reflections whose
axes lie \emph{along} edge directions exchange them; the even rotations and the
three reflections whose axes lie at $30^\circ$ to the edges do not.
\end{corollary}

\begin{proof}
Immediate from Theorem~\ref{thm:star}(iii), since the labelling is a function of
triangle orientation and $g$ maps $\Star(v)$ to $\Star(g v)$. The classification
of which $g$ exchange orientations: writing a reflection as $d\mapsto m-d$ on
direction indices, it fixes a direction iff $2d\equiv m \pmod 6$, which is
solvable iff $m$ is even; the $m$ even reflections (axis along an edge
direction) exchange up and down triangles and the $m$ odd ones do not
(verified on $\mathrm{up}(0,0)$ for all twelve group elements).
\end{proof}

\begin{remark}
Corollary~\ref{cor:dual} is the structural content of the paper's two hexagonal
theorems. It is \emph{not} that $\hex$ and $\hexR$ are two boundaries of one
model; they are one tile in the two orientation conventions. Consistently, the
paper \emph{defines} $\tilde d^\nu_{\la\mu} := d^{\nu'}_{\la'\mu'}$ on transposed
partitions (\texttt{ktiles.tex}, lines 757--760), and the rectangle bookkeeping
matches exactly: with $k^\ast=n-k$ and $n^\ast=n$, the $\hexR$ ambient rectangle
$(k-1)\times(n-k)$ transposes to $(n-k)\times(k-1) =
k^\ast\times(n^\ast-k^\ast-1)$, the $\hex$ ambient rectangle. So
Theorem~\texttt{hexR-tiling} is Theorem~\texttt{hex-tiling} conjugated by
$k\mapsto n-k$, i.e.\ by the Grassmannian duality
$\mathrm{Gr}(k,n)\cong\mathrm{Gr}(n-k,n)$.
\end{remark}

\section{Protrusion is a half-overlap, and the trapezoid is a clipping}

\begin{theorem}[Protrusion]\label{thm:protrude}
Let $W$ be a lattice line and let $H^{+},H^{-}$ be the two closed half-planes it
bounds. For a lattice vertex $v$:
\[
 \Star(v)\subseteq H^{+} \iff v\in H^{+}\setminus W, \qquad
 \Star(v)\subseteq H^{-} \iff v\in H^{-}\setminus W,
\]
and $\Star(v)$ meets both open half-planes if and only if $v\in W$. In that
case exactly three of its six triangles lie in each half-plane, and
$W\cap\Star(v)$ is the segment $[v-e,v+e]$ of length $2$, where $e$ is the unit
vector along $W$.
\end{theorem}

\begin{proof}
Choose the basis so that $W=\{a=0\}$ and $e=e_1$ (possible since $W$ is a
lattice line, hence parallel to some $D_j$; take $e_1=D_j$). The six neighbours
$v+D_j$ of $v=(a_v,b_v)$ have $a$-coordinates $a_v+(1,0,-1,-1,0,1)$, so
\[
 \Star(v)\subseteq\{\,a_v-1\le a\le a_v+1\,\},
\]
with both extremes attained. Hence $\Star(v)\subseteq\{a\ge0\}$ iff $a_v\ge1$,
and $\Star(v)\subseteq\{a\le0\}$ iff $a_v\le-1$; the remaining case is $a_v=0$.
If $a_v=0$ the six triangles $T_j$ split by the sign of the $a$-coordinates of
their non-$v$ vertices: three have an $a=-1$ vertex and three an $a=+1$ vertex,
and no triangle has both (consecutive $D_j,D_{j+1}$ never have $a$-coordinates
$+1$ and $-1$). So the split is $3+3$. Finally $W\cap\Star(v)$ contains $v$ and
the two neighbours $v\pm e_1$ (the only neighbours on $W$), and is convex, so it
is $[v-e_1,v+e_1]$.
\end{proof}

\begin{corollary}[why ``slightly smaller'' is enough]
Because the obstruction occurs only for $v\in W$, a hexagon can protrude past
the wall by exactly one lattice step and no more. Enlarging the ambient
rectangle by one in the direction normal to the wall moves every admissible
centre to $a_v\ge1$ and removes the obstruction. This is the paper's
Theorem~\texttt{hex-tiling} hypothesis and the footnote ``For technical
reasons, we require partitions to be slightly smaller.''
\end{corollary}

\begin{theorem}[The trapezoid is a truncation]\label{thm:trunc}
With $W=\{a=0\}$ the wall and the puzzle in $\{a\ge0\}$, let $v\in W$ so that
$\Star(v)$ straddles. Then the \S\texttt{ss:protrude} trapezoid is exactly
\[
 \trap \;=\; \hex \cap \{a\ge 0\},
\]
as labelled tiles: the three retained boundary edges keep their labels
unchanged, the three edges lying in $\{a<0\}$ are discarded, and the only new
boundary edges are the two unit edges of $W\cap\Star(v)=[v-e_1,v+e_1]$, carrying
label $0$ on $[v,v+e_1]$ and label $1$ on $[v-e_1,v]$.
\end{theorem}

\begin{proof}
Verified exactly, by comparing covered unit triangles and labelled boundary
edges. With $\hex$ anchored so that $v=(0,1)$: $\Star(v)$ covers the six
triangles $\{U(0,0),U(0,1),U(-1,1),D(0,0),D(-1,0),D(-1,1)\}$, of which exactly
$\{U(0,0),U(0,1),D(0,0)\}$ lie in $\{a\ge0\}$; the trapezoid covers precisely
these three. The retained boundary edges are
$[(1,0),(1,1)]$ with label $0$, $[(1,1),(0,2)]$ with label $1$, and
$[(0,0),(1,0)]$ with label $1$ --- identical in $\hex$ and in $\trap$. The
discarded edges are $[(-1,2),\cdot]$ label $0$, $[(-1,1),\cdot]$ label $1$,
$[(0,0),\cdot]$ label $0$. The new edges are $[(0,1),(0,2)]$ label $0$ and
$[(0,0),(0,1)]$ label $1$, which are the two halves of $W\cap\Star(v)$ above and
below $v=(0,1)$. The area accounting is consistent:
$\#\trap = 3 = \tfrac12\#\hex$, as Theorem~\ref{thm:protrude} requires.
\end{proof}

\begin{remark}[one bit of new data, and no freedom in it]\label{rem:onebit}
Clipping does introduce something not present in $\hex$: the two cut edges were
\emph{interior} edges of the hexagon (the radii $[v,v\pm e_1]$), hence
unlabelled, and become boundary edges requiring labels. Theorem~\ref{thm:trunc}
records that these are $0$ above $v$ and $1$ below. This is the trapezoid's
entire content beyond the hexagon --- one bit --- and it is not a free
parameter: the paper lists a single trapezoid, and \S\ref{sec:protrude} verifies
computationally that this single choice reproduces $\dd$.
\end{remark}

\section{There is no reflection to reflect in}\label{sec:refl}

A K-matrix is a boundary vertex that \emph{reflects} bulk lines and satisfies the
reflection equation
$K_2 R_{21} K_1 R_{12} = R_{21} K_1 R_{12} K_2$. Pre-registered prediction P2
said the trapezoid plays this role. It does not, for three independent reasons.

\begin{proposition}[Zero moduli]\label{prop:moduli}
The boundary piece is determined by the bulk piece and the wall:
$\trap = \hex\cap H^{+}$ (Theorem~\ref{thm:trunc}), and the clipping is unique
once the wall is fixed (Theorem~\ref{thm:protrude} gives a unique straddling
configuration, with a unique cut). A K-matrix, by contrast, is an
\emph{independent} solution of the reflection equation: it carries its own
spectral parameter and generically belongs to a positive-dimensional family for
a given $R$. The PY boundary piece has no spectral parameter and no family.
\end{proposition}

\begin{proposition}[The bulk tile set is chiral]\label{prop:chiral}
No reflection of the triangular lattice preserves the KTW tile set.
\end{proposition}

\begin{proof}
The KTW rhombus rule reads: edges are labelled ``$0$ if clockwise of an acute
angle and $1$ if clockwise of an obtuse angle'' (\texttt{ktiles.tex}, line 577).
A reflection reverses orientation, hence exchanges ``clockwise'' for
``counterclockwise'', and therefore sends each rhombus to the rhombus on the
same pair of direction classes carrying the \emph{complementary} labelling. None
of the three complementary rhombi is a KTW piece. Verified for all six
reflections: each sends all three rhombi outside the set (while permuting the
$0$- and $1$-triangles among themselves). The six rotations, by contrast, all
preserve the set.
\end{proof}

\begin{proposition}[The hexagon is not wall-reflection-invariant]\label{prop:notinv}
Let $\sigma_W$ be the reflection in the wall $W$. As a \emph{region},
$\Star(v)$ is $\sigma_W$-invariant for $v\in W$. As a \emph{labelled tile} it is
not: $\sigma_W(\hex)=\hexR$.
\end{proposition}

\begin{proof}
$W$ is a lattice line, so $\sigma_W$ has its axis along an edge direction;
by Corollary~\ref{cor:dual} such a reflection exchanges the triangle
orientations, hence exchanges $\hex$ and $\hexR$. Regionally $\sigma_W$ fixes
$v$ and permutes its neighbours, so it fixes $\Star(v)$.
\end{proof}

\begin{remark}[what Proposition~\ref{prop:notinv} costs the reflection reading]
This is the sharpest of the three. Under the wall reflection the hexagon does
not return to itself but to $\hexR$ --- a \emph{different} tile, which computes a
\emph{different} coefficient ($\tilde d$ rather than $d$) in a
\emph{different} ambient rectangle. So the two halves of a straddling hexagon
are mirror images as shapes but carry mismatched labels under the mirror: the
clipping is not a fold, and there is no reflection symmetry at the wall for a
reflection equation to express.
\end{remark}

\begin{proposition}[The label $2$ is an indicator, not an amplitude]
In the \S\texttt{ss:protrude} tile set, a $1$-edge in the wall's direction class
$\{1,4\}$ carries label $1$ if and only if it lies on $W$, and label $2$ if and
only if it is interior.
\end{proposition}

\begin{proof}
Counterclockwise traversal puts a tile's interior on the left of each directed
edge. Hence a tile presents a direction-$1$ edge exactly when it lies to that
edge's upper left, and a direction-$4$ edge exactly when it lies to its lower
right. The paper relabels every $1$-edge on a tile's lower-right side to $2$
(so direction-$1$ $1$-edges are always $2$: pieces \verb|\bb1\b3\b5|,
\verb|\bb1\r2\b4\r5|, \verb|\bb1\r2\bb4\r5|), and duplicates every piece with a
$1$-edge on its upper-left side (direction $4$) into a label-$1$ and a label-$2$
version. Now let $e$ be an interior lattice edge in this class: the tile above
left of $e$ presents it in direction $1$, hence labels it $2$, so the tile below
right must also label it $2$, i.e.\ must be the label-$2$ version. If instead
$e\subseteq W$ there is no tile above left, and the tile below right labels it
$1$ --- consistent with Convention~\ref{conv:bdry}, under which the wall carries
$\sigma(\la)\in\{0,1\}^n$. So label $1$ on this class occurs exactly on $W$.
In particular $\trap$, whose direction-$4$ $1$-edge carries label~$1$ with no
label-$2$ counterpart, can be placed only at the wall --- the paper's ``the new
trapezoid piece cannot be used except at the left boundary''. There is no sum
over boundary states anywhere in this: the label is a deterministic function of
wall-adjacency.
\end{proof}

\section{There is no vertex model on these tiles}\label{sec:nomodel}

Pre-registered prediction P1 said the bulk set does not satisfy a plain
Yang--Baxter equation \emph{for any assignment of spectral parameters and
weights}. That prediction is not well posed: the tiles are faces of a tiling
model, all of weight $1$, and the paper's theorems count cardinalities; there is
no $R$ to test until one fixes a map from tiles to a matrix on a tensor
product, and the paper supplies none. (This is exactly the failure the brief
warned about --- check that the objects exist before grading the chances.)

Here is the well-posed replacement. In a vertex model the edge states are line
occupations, and a line entering a vertex must leave it, so each vertex carries
an \emph{even} number of occupied edges in each direction class. Taking the
$1$-edges as the lines and the classes $H=\{0,3\}$, $N=\{1,4\}$,
$W'=\{2,5\}$:

\begin{proposition}\label{prop:noconserve}
Ten of the seventeen tiles violate line conservation. The $0$-triangles and all
rhombi conserve (counts $(0,0,0)$ and one class $=2$); the $1$-triangles,
$\Ktile$, $\tritile$, $\hex$, $\hexR$, $\mathsf{hex}_2$ and $\trap$ each carry
exactly \emph{one} $1$-edge in each of the three classes, i.e.\ $(1,1,1)$ --- odd
in all three.
\end{proposition}

\begin{proof}
Direct count from the step sequences of \S\ref{sec:objects}.
\end{proof}

So the $1$-edges form a graph in which rhombi are straight segments and every
other nontrivial tile is a $Y$-junction: a branching structure, not a conserving
one. The tile set is therefore not itself a vertex model, and the question
``does it satisfy YBE'' has no object. This is consistent with the paper's
measured silence (the brief's census: \emph{Yang--Baxter}, \emph{R-matrix},
\emph{integrable}, \emph{vertex model}, \emph{transfer matrix},
\emph{lattice model} --- $0$ occurrences each across $1790$ lines, against six
positive controls firing on the same instrument).

\begin{remark}[scope --- what this does \emph{not} say]
It does not say that Grassmannian puzzles are non-integrable. They are known to
be integrable, and the paper's own $\tritile$ piece is due to
Wheeler--Zinn-Justin, who obtained it \emph{from} an integrable lattice model.
Proposition~\ref{prop:noconserve} says only that the integrable structure does
not live on \emph{these faces with these labels as line occupations}; any
integrable formulation must pass through a different presentation, and
Pylyavskyy--Yang supply none for the two new pieces. Whether the hexagon has an
$R$-matrix provenance in some other presentation is open and is \emph{not}
settled here.
\end{remark}

\section{The \S\texttt{ss:protrude} alternative, verified}\label{sec:protrude}

The paper offers its second fix without a theorem or a proof: ``Here is the
complete set of tiles'', followed by a figure of twelve pieces and three
sentences of explanation. The claim is that with the refined alphabet
$\{0,1,2\}$ one may use the \emph{natural} side length rather than enlarging it.
I tested it.

The twelve pieces are: the two $0$-triangles; the $1$-triangles
\verb|\b0\b2\b4|, \verb|\b0\b2\bb4|, \verb|\bb1\b3\b5|; the rhombi
\verb|\b0\r1\b3\r4|, \verb|\b2\r3\b5\r0|, \verb|\bb1\r2\b4\r5|,
\verb|\bb1\r2\bb4\r5|; and the green pieces
\verb|\r1\b2\r3\b4\r5\b0| $(=\hex)$, \verb|\r1\b2\r3\bb4\r5\b0|,
\verb|\r1\b2\r4\b4\b0| $(=\trap)$.

\begin{proposition}[computed]\label{prop:protrude}
For all $\la,\mu,\nu$ inside the ambient rectangle $k\times m$, with
$|\nu|\le|\la|+|\mu|$, the number of refined-set tilings at side $n=k+m$ equals
the number of $\hex$-tilings at side $n+1$ equals $\dd$.
\emph{Census domain:} $(k,m)=(2,2)$, all $187$ admissible triples, $44$ of them
with $\dd\neq0$; and $(k,m)=(2,3)$, all $875$ admissible triples, $143$ of them
with $\dd\neq0$. Agreement $187/187$ and $875/875$ for both the refined set and
the enlarged $\hex$ set, against Engine~B.
\end{proposition}

I also ran the paper's own \S\texttt{ss:protrude} example,
$\la=(2,1)$, $\mu=(4,2)$, $\nu=(4,3,1)$, $k=3$, which is where the protrusion is
exhibited:

\begin{center}\small
\begin{tabular}{lr}
\hline
$\dd$ from Engine~B (independent) & $4$\\
$\hex$-puzzles at side $7$ (the ``natural'' side; theorem's hypothesis fails, $\mu_1=4>3$) & $2$\\
$\hex$-puzzles at side $8$ (the theorem's prescribed side) & $4$\\
refined \S\texttt{ss:protrude} set at side $7$ & $4$\\
\hline
\end{tabular}
\end{center}

\noindent
So at the natural side length the plain $\hex$ set loses exactly two tilings ---
the protrusion, quantified --- and both of the paper's fixes recover all four.

\begin{corollary}[the shrink is not intrinsic]\label{cor:notintrinsic}
The one-unit shrink of the ambient rectangle is not a structural feature of the
$\hex$ rule. It is one of two equivalent bookkeeping devices for the same
clipping: move the wall away from the hexagon's centre, or clip the hexagon and
mark wall-adjacency in the alphabet.
\end{corollary}

\section{The pre-registered predictions, graded}

\begin{description}
\item[P1] \emph{The bulk set does not satisfy a plain YBE for any assignment of
spectral parameters and weights. Predicted: confirmed.} \textbf{SPLIT.} The
conclusion stands but the reason in the brief was wrong: there is no weight
object to quantify over, so the statement as written is ill-posed. The
well-posed replacement (Proposition~\ref{prop:noconserve}) is stronger in one
direction --- there is no vertex model on these tiles at all --- and weaker in
another: it says nothing about other presentations.

\item[P2] \emph{With the trapezoid as $K$ and the rhombi as $R$, the set
satisfies a reflection equation. Predicted: confirmed; the risky arm.}
\textbf{REFUTED}, on three independent grounds: zero moduli
(Proposition~\ref{prop:moduli}), chirality of the bulk
(Proposition~\ref{prop:chiral}), and $\sigma_W(\hex)=\hexR$
(Proposition~\ref{prop:notinv}). This is the arm I most wanted to be wrong
about, and it died.

\item[P3] \emph{The one-unit shrink is the boundary site: $\hex$ and $\hexR$ are
the two one-sided boundaries of the same bulk model. Predicted: confirmed.}
\textbf{REFUTED as stated.} $\hex$ and $\hexR$ are one tile in the two triangle
orientations (Corollary~\ref{cor:dual}), and the second theorem is the first
conjugated by $k\mapsto n-k$ with transposed partitions. There is one boundary
phenomenon seen twice through duality, not two boundaries. What survives: the
protrusion really is one-sided, and Theorem~\ref{thm:protrude} says why. What
additionally dies: the shrink is not intrinsic
(Corollary~\ref{cor:notintrinsic}).

\item[P4] \emph{The degree-inequality flip is the $K$-homology/$K$-theory
duality and not anything about the boundary --- i.e.\ P4 predicts this clue is a
red herring for Q569. Predicted: confirmed.} \textbf{CONFIRMED.} The flip is
product versus coproduct: $|\nu|\ge|\la|+|\mu|$ for $G_\la G_\mu$ and
$|\nu|\le|\la|+|\mu|$ for $\Delta(G_\nu)$ (\texttt{ktiles.tex}, lines 725--728),
both reproduced in Engine~B. It is orthogonal to the wall.
\end{description}

\noindent\textbf{Pre-committed falsifiers.} ``If the bulk tiles do satisfy a
plain YBE'': not reached, since no YBE object exists (P1 SPLIT).
``If no consistent weight assignment exists at all'': this is what happened, and
I report it as Proposition~\ref{prop:noconserve} rather than as a refutation of
the paper, which claims nothing about weights. ``If P2 holds but only after the
trapezoid is allowed in the bulk'': not reached; P2 failed outright.

\section{Gaps, precisely stated}

\begin{enumerate}
\item \textbf{Proposition~\ref{prop:protrude} is \texttt{computed}, not
\texttt{proved}.} I verified the \S\texttt{ss:protrude} equivalence on two
ambient rectangles ($2\times2$, $2\times3$) and on the paper's own example. I do
not have a bijection. The natural candidate --- ``clip every hexagon whose
centre lies on the wall, and read the extra leading $0$ of
$\sigma(\cdot,n+1,k)$ as a frozen strip of $0$-triangles'' --- I checked only to
the extent that the corner triangle $U(0,0)$ of the enlarged puzzle is forced to
be a $0$-triangle (the only tile admitting two $0$-labelled boundary edges
there). Propagating that to a full frozen region, and matching the clipped
hexagons across the two models, is not done.

\item \textbf{The integrability question is settled only in the presentation I
tested.} Proposition~\ref{prop:noconserve} rules out a line-conserving vertex
model on these faces with these labels. It does not rule out an $R$-matrix for
$\hex$ in another presentation, e.g.\ the one Wheeler--Zinn-Justin use for
$\tritile$. I did not read \cite{WZ}; I have not opened it, and no claim here
depends on it.

\item \textbf{Proposition~\ref{prop:moduli} contains a judgement, not a
theorem.} ``A K-matrix generically belongs to a positive-dimensional family'' is
standard but I state it from general knowledge, not from a source read this
session. The hard part --- that $\trap$ is determined by $\hex$ and $W$ --- is
Theorem~\ref{thm:trunc} and is proved.

\item \textbf{Census domains are small.} The four-theorem control reaches
$n\le6$, $k\le3$; the \S\texttt{ss:protrude} check reaches ambient $2\times3$.
A coincidence surviving to $k=3$ but failing at $k=4$ would not be visible here.
I report the domain rather than ``small cases''.
\end{enumerate}

\section{Instrument faults found this session}

\begin{enumerate}
\item \textbf{A reflection implemented as a step-order reversal.} My group
action reversed a tile's step sequence without negating the directions
($d\mapsto d+3$). Reversing the order alone permutes a closed walk's edges into
a \emph{different} polygon --- which still closes and is still
counterclockwise, so neither my closure assertion nor my signed-area assertion
detected it. The wrong code reported that the KTW set is reflection-closed for
no $g$ but that $\sigma_W(\hex)=\hex$; the correct code reports
$\sigma_W(\hex)=\hexR$, which reverses the sign of the central argument in
\S\ref{sec:refl}. The lesson in the form I want to keep it:
\emph{two invariants of a polygon that are both invariant under edge
permutation cannot detect an edge permutation} --- the control must be a
property that distinguishes the orbit, here
$\mathrm{act}(\mathrm{act}(x,g),g^{-1})=x$ (now checked on all $204$
(tile, $g$) pairs) plus an independent orientation test
(does $g$ send an up-triangle to an up-triangle?).

\item \textbf{A correct answer reached through a wrong intermediate.} Before
fixing (1) I had written that the hexagon is wall-reflection-\emph{invariant}
and that this is ``why clipping it in half is consistent''. The clipping is
consistent, and $\trap=\hex\cap H^{+}$ is true, but not for that reason. The
true reason is Theorem~\ref{thm:protrude}: the wall passes through the centre,
so the region splits $3+3$. Region-invariance and label-invariance are different
statements and I had merged them.

\item \textbf{A relative-path emitter default.} \texttt{tracecheck.emit.init}
defaults its log directory to a \emph{relative} path, so my first call --- made
with cwd $=$ \texttt{projects/} --- wrote the session trajectory to
\texttt{projects/state/trajectory/} instead of \texttt{state/trajectory/}. The
file-listing check I ran immediately afterwards looked in the intended
directory and showed only older sessions, which is what caught it. Re-pinned
with an absolute \texttt{log\_dir} and a stable session tag.

\item \textbf{A float where an integer was expected, harmlessly.} Engine~B
returns $\pm1$ sign factors as Python floats whenever the exponent
$|\nu|-|\la|-|\mu|$ is negative (which is always, for $d$). The comparisons
$4.0 == 4$ behave correctly, so no result is affected, but the printed values
read ``$4.0$'' and a future reader should not take that as a non-integral
structure constant.
\end{enumerate}

\section{What I would do next}

The bijection in Gap 1 is the obvious next target and looks tractable: by
Theorem~\ref{thm:protrude} the only configurations differing between the two
models are hexagons centred on the wall, and by Theorem~\ref{thm:trunc} the
clipping is label-preserving, so the content is entirely in showing that the
enlarged puzzle's extra column is frozen. If that works it converts
Proposition~\ref{prop:protrude} from \texttt{computed} to \texttt{proved} and
supplies the proof the paper omits.

The second target is the one I deliberately did not reach for: whether $\hex$
has an $R$-matrix provenance in the Wheeler--Zinn-Justin presentation rather
than on its own faces. Theorem~\ref{thm:star} is encouraging for this, because
the vertex star is the natural face of the \emph{dual} (honeycomb) lattice ---
and honeycomb is where Grassmannian puzzle integrability usually lives. That is
a conjecture and is recorded as one.

\begin{thebibliography}{9}
\bibitem{PY} P.~Pylyavskyy and J.~Yang, \emph{Puzzles in $K$-homology of
Grassmannians}, arXiv:1801.07667. E-print read in full this session
(\texttt{ktiles.tex}, \texttt{hex-tiling.tex}); all tile data, theorem
statements, conventions and quotations above are from it.
\bibitem{Buch} A.~S.~Buch, \emph{A Littlewood--Richardson rule for the
$K$-theory of Grassmannians}. Cited only for the definition of $G_\la$ via
set-valued tableaux, which I took from \cite{PY} lines 711--716 rather than from
Buch directly; \emph{not read this session}.
\bibitem{WZ} M.~Wheeler and P.~Zinn-Justin, cited by \cite{PY} for the
$\tritile$ piece. \emph{Not read this session}; no claim here depends on it.
\bibitem{KTW} A.~Knutson, T.~Tao and C.~Woodward, cited by \cite{PY} for the
original puzzle pieces, which I extracted from \cite{PY} figure LR-tiles.
\emph{Not read this session.}
\end{thebibliography}

\end{document}
